% ------------------------------------------------------------
% RISC-V instruction encoding diagrams
% ------------------------------------------------------------

\usepackage{tikz}

% Width of one bit in the current encoding diagram.
% The complete instruction occupies \linewidth.
\newlength{\riscvbitwidth}

% Current position, measured from the MSB.
\newcount\riscvencodingpos

% Number of bits emitted so far.
\newcount\riscvencodingbits

% Width of the current instruction.
\newcount\riscvinstructionbits


% ------------------------------------------------------------
% TikZ styles
% ------------------------------------------------------------

\tikzset{
  riscv/encoding/.style={
    x=\riscvbitwidth,
    y=1cm,
    font=\sffamily,
    line cap=butt,
    line join=miter
  }
}

\tikzset{
  riscv/field/.style={
    draw=ephemink,
    line width=0.35pt
  }
}

\tikzset{
  riscv/value/.style={
    font=\ttfamily\small,
    inner sep=0pt
  }
}

\tikzset{
  riscv/bits/.style={
    font=\sffamily\scriptsize,
    inner sep=0pt
  }
}

\tikzset{
  riscv/format/.style={
    font=\sffamily\small,
    anchor=west,
    inner sep=0pt
  }
}


% ------------------------------------------------------------
% RISC-V encoding environment
%
% Usage:
%
%   \begin{riscvencoding}[R-type]{MULH}{32}
%     \field{0000001}{7}
%     \field{rs2}{5}
%     \field{rs1}{5}
%     \field{001}{3}
%     \field{rd}{5}
%     \field{0110011}{7}
%   \end{riscvencoding}
%
%   \begin{riscvencoding}[C-type]{C.ADDI}{16}
%     \field{01}{2}
%     \field{000}{3}
%     \field{imm[5]}{1}
%     \field{rs1}{5}
%     \field{imm[4:0]}{5}
%   \end{riscvencoding}
%
% Arguments:
%
%   #1 = optional instruction format
%   #2 = instruction name
%   #3 = instruction width in bits
%
% The complete diagram occupies \linewidth.
% ------------------------------------------------------------

\newenvironment{riscvencoding}[3][]{%

  % Save instruction width.
  \riscvinstructionbits=#3\relax

  % One bit occupies an equal fraction of the available width.
  \setlength{\riscvbitwidth}{%
    \dimexpr\linewidth/#3\relax
  }%

  % Start at the most-significant bit.
  \riscvencodingpos=0%

  % No fields emitted yet.
  \riscvencodingbits=0%

  % Save format and instruction name.
  \def\riscvformat{#1}%
  \def\riscvinstruction{#2}%

  \begin{center}
  \begin{tikzpicture}[riscv/encoding]

}{%

  % Verify that the fields occupy the complete instruction.
  \ifnum\riscvencodingbits=\riscvinstructionbits
  \else
    \PackageError{riscvencoding}
      {Instruction encoding contains \the\riscvencodingbits\space bits}
      {The encoding must contain exactly
       \the\riscvinstructionbits\space bits.}%
  \fi

  % Optional instruction format and name.
  \ifx\riscvformat\empty
  \else
    \node[
      riscv/format
    ] at (
      \number\numexpr\riscvinstructionbits+1\relax,
      0.52
    ) {\riscvformat};

    \node[
      riscv/format
    ] at (
      \number\numexpr\riscvinstructionbits+1\relax,
      0.16
    ) {\instr{\riscvinstruction}};
  \fi

  \end{tikzpicture}
  \end{center}
}


% ------------------------------------------------------------
% \field{label}{width}
%
% #1 = field label / encoded value
% #2 = field width in bits
%
% Examples:
%
%   \field{0000001}{7}
%   \field{rs2}{5}
%   \field{imm[4:0]}{5}
%
% The label is rendered exactly as supplied.
% ------------------------------------------------------------

\newcommand{\field}[2]{%

  % ----------------------------------------------------------
  % Calculate field boundaries
  % ----------------------------------------------------------

  % Position immediately after this field.
  \pgfmathtruncatemacro{\riscvendpos}{%
    \riscvencodingpos + #2%
  }%

  % Most-significant bit of this field.
  \pgfmathtruncatemacro{\riscvstartbit}{%
    \riscvinstructionbits - 1 - \riscvencodingpos
  }%

  % Least-significant bit of this field.
  \pgfmathtruncatemacro{\riscvendbit}{%
    \riscvinstructionbits - \riscvendpos
  }%

  % ----------------------------------------------------------
  % Validate field
  % ----------------------------------------------------------

  \ifnum\riscvendpos>\riscvinstructionbits
    \PackageError{riscvencoding}
      {Encoding field exceeds instruction width}
      {A field extends beyond the
       \the\riscvinstructionbits-bit instruction.}%
  \fi

  % ----------------------------------------------------------
  % Field rectangle
  % ----------------------------------------------------------

  \draw[
    riscv/field
  ]
    (\riscvencodingpos,0)
    rectangle
    (\riscvendpos,0.72);

  % ----------------------------------------------------------
  % Field label
  %
  % The field is 0.72 units high, so 0.36 is its vertical
  % center.
  % ----------------------------------------------------------

  \node[
    riscv/value
  ] at (
    \riscvencodingpos + #2/2,
    0.36
  ) {#1};

  % ----------------------------------------------------------
  % Field bit numbers
  %
  % Both numbers are placed inside the field, slightly inset
  % from its boundaries.
  %
  % Example:
  %
  %   31        25 24        20
  %   ┌───────────┬────────────┐
  %
  % The left number uses south-west anchoring and the right
  % number uses south-east anchoring, so adjacent numbers
  % remain separated.
  % ----------------------------------------------------------

  % Most-significant bit.
  \node[
    riscv/bits,
    anchor=south west,
    xshift=2pt
  ] at (
    \riscvencodingpos,
    0.98
  ) {\riscvstartbit};

  % Least-significant bit.
  \node[
    riscv/bits,
    anchor=south east,
    xshift=-2pt
  ] at (
    \riscvendpos,
    0.98
  ) {\riscvendbit};

  % ----------------------------------------------------------
  % Advance to next field
  % ----------------------------------------------------------

  \riscvencodingpos=\riscvendpos

  \advance\riscvencodingbits by #2\relax
}
